\subsubsection{Harsanyi-Selten Nash equilibria as a benchmark}

The second benchmark is symmetric Nash equilibria.
Unlike the cognitive hierarchy model, this solution concept does not require any previous data to generate predictions.
It is also flexible: symmetric Nash equilibria exist for any symmetric two-player game with a finite number of actions \citep{nash1951noncooperative}.
They can be computed systematically across all game types in our dataset.
Furthermore, in our setting, all games are played independently by participants (AI and human), making symmetry a natural assumption that suggests a ``consistent common belief'' across the population \citep{STAHL1994}.

Since many games have multiple symmetric Nash equilibria, we require a systematic method to select a single equilibrium prediction. 
Indeed, one game in $S$ has $\MaxSymmNash$ symmetric equilibria. 
Unfortunately, there is no universally agreed-upon criterion for selecting the ``optimal'' symmetric equilibrium across all games \citep{camerer2003behavioral, tadelis2013game}.

We employ a slightly modified version of the equilibrium selection procedure developed by \citet{HarsanyiSelten1988}, which provides a principled approach grounded in stability and focal-point considerations. 
It guarantees the selection of a unique Nash equilibrium for every game in our dataset and can be slightly modified to ensure symmetry. 
The procedure also prioritizes equilibria that are stable in the sense of \citet{schelling1960strategy}, specifically favoring equilibria that are either payoff dominant (maximizing joint welfare) or risk dominant (robust to strategic uncertainty). 
This is appealing because many of the games---particularly those with bonus rules ``Equal'' or ``Coord. Low''---often have clear equilibria that are both payoff and risk dominant, which humans often choose \citep{camerer2003behavioral}. 

The Harsanyi-Selten procedure operates through a multi-stage filtering process, progressively narrowing the set of candidate equilibria. 
The procedure first identifies all Nash equilibria, then applies filters based on Pareto efficiency, symmetry requirements, and risk dominance, and finally employs tracing methods to resolve any remaining ties. 
See the pseudocode in Appendix~\ref{app:nash-selector} for details on the implementation.

We first use open-source software \citep{gambit2025} to calculate all Nash equilibria for the \NGamesS{} games in set $S$. 
We then apply the Harsanyi-Selten procedure to each game's set of equilibria, producing a single equilibrium distribution for \GamesConverged{} of the games.\footnote{
In fewer than 1\% of games, degeneracy issues prevented the code from converging.
Following our preregistration, we discard these games in our analysis when comparing the equilibria to the strategic \aisub{s}.}
Of these \GamesConverged{} games, \CountUniqueEq{} have unique symmetric equilibria.
The selection procedure was unnecessary in these cases. 
Among the remaining games with multiple symmetric equilibria, the procedure selects payoff-dominant equilibria---those Pareto superior to all alternatives---in \CountPayoffDomEq{} games, and risk-dominant equilibria in \CountRiskDomEq{} games. 
Finally, $\PropSelectedPure$\% of the equilibria selected by the Harsanyi-Selten procedure have pure strategies, with the remainder employing mixed strategies.

Similar to the predictions from the cognitive hierarchy model, the equilibrium distributions are highly diverse across games.
This can be seen in Figure~\ref{fig:eq_plot} in the appendix, which is of the same format as Figure~\ref{fig:ch_plot}.
